

\begin{frame}
  \frametitle{TL;DR Simplification of the \gls{bte}}
  By inserting const-value properties, noting the division of $4\pi$ for the scattering cross section to keep units consistent, we are left with:
  \begin{equation} \label{eq:fully-simplified-bte}
    \ohat \cdot \grad \psisimple + \Sigma_t \left(\rvec\right)\psisimple =\underbrace{\frac{\Sigma_s\left(\rvec\right)}{4\pi}\phi\left(\rvec\right) + \frac{Q\left(\rvec\right)}{4\pi}}_{\equiv S\left(\rvec\right)}
  \end{equation}
  where:
  \begin{equation}
    \label{eq:scalar-flux}
    \phi\left(\rvec\right)\equiv \int_{4\pi}\psi \left(\rvec,\ohat\right) \dd \ohat
  \end{equation}
  \Cref{eq:scalar-flux} is know as the \textit{angle-integrated} (or \textit{scalar}) \textit{flux}.
\end{frame}